\documentclass[12pt,a4paper]{article}

% ---- Encoding & Font ----
\usepackage{iftex}
\ifPDFTeX
  \usepackage[utf8]{inputenc}
  \usepackage[T1]{fontenc}
\else
  \usepackage{fontspec}  % XeTeX/LuaTeX (tectonic): Unicode nativo (·, —, ¿, ª, §)
\fi
\usepackage[spanish]{babel}
\usepackage{csquotes}

% ---- Page layout ----
\usepackage[top=2.5cm, bottom=2.5cm, left=2.5cm, right=2.5cm]{geometry}
\usepackage{setspace}
\onehalfspacing

% ---- Math ----
\usepackage{amsmath, amssymb, amsthm}
\usepackage{mathtools}
\usepackage{physics}

% ---- Graphics ----
\usepackage{graphicx}
\usepackage{tikz}
\usepackage{float}
\usepackage{caption}

% ---- Tables ----
\usepackage{array}
\usepackage{booktabs}
\usepackage{tabularx}
\usepackage{colortbl}

% ---- Colours ----
\usepackage{xcolor}
\definecolor{notebg}{HTML}{FFF3CD}
\definecolor{noteborder}{HTML}{FFC107}
\definecolor{boxred}{HTML}{F8D7DA}
\definecolor{boxredborder}{HTML}{F5C6CB}

% ---- Boxes ----
\usepackage{mdframed}
\usepackage{tcolorbox}
\tcbuselibrary{most}

\newtcolorbox{keybox}[1][]{
  colback=blue!5,
  colframe=blue!70!black,
  arc=4pt,
  boxrule=1pt,
  fonttitle=\bfseries,
  title={#1}
}

\newtcolorbox{warnbox}[1][]{
  colback=boxred,
  colframe=boxredborder,
  coltitle=black!75,
  arc=4pt,
  boxrule=1pt,
  fonttitle=\bfseries,
  title={#1}
}

\newtcolorbox{notebox}[1][]{
  colback=notebg,
  colframe=noteborder,
  coltitle=black!75,
  arc=4pt,
  boxrule=1pt,
  fonttitle=\bfseries,
  title={#1}
}

% ---- Hyperlinks & TOC ----
\usepackage[hidelinks]{hyperref}
\usepackage{bookmark}

% ---- Enumerate ----
\usepackage{enumitem}

% ---- Miscellaneous ----
\usepackage{icomma}
\usepackage{siunitx}
\usepackage{textcomp}
\usepackage[bottom]{footmisc}

% ---- Title ----
\title{\textbf{Clase 7: Normas de Matrices} \\[0.3em]
        \large{Normas operador, la fórmula para $\|\cdot\|_\infty$, y error vs. residuo}}
\author{Transcripto y expandido --- Curso Métodos Numéricos (OpenFING)}
\date{}

\begin{document}

\maketitle
\thispagestyle{empty}
\newpage

\tableofcontents
\newpage

\section{Cierre de matrices dispersas y tridiagonales}

Se define informalmente la \textbf{densidad} de $A \in \mathbb{R}^{n\times
n}$ como el cociente entre su cantidad de entradas no nulas y $n^2$
(\textbf{dispersidad} = $1$ menos eso); una matriz es \textbf{dispersa}
cuando su dispersidad es cercana a $1$. Para esas matrices conviene
almacenar sólo las ternas $(i,j,a_{ij})$ no nulas.

\begin{notebox}[Matrices de banda / tridiagonales]
$A$ es tridiagonal si $a_{ij}=0$ para todo $|i-j|>1$. Si se le avisa al
código esta estructura, la eliminación gaussiana con pivoteo parcial se hace
en \textbf{tiempo lineal} en vez de $O(n^3)$: cada paso pivotea comparando
sólo con la fila de abajo, y las actualizaciones son de costo constante por
fila. \textbf{Si no se avisa la estructura, el algoritmo cae en $O(n^3)$
igual}, aunque la matriz sea tridiagonal.
\end{notebox}

\section{Normas vectoriales: repaso de axiomas}

Una \textbf{norma} en un espacio vectorial $V$ sobre $\mathbb{R}$ es
$\|\cdot\|: V \to \mathbb{R}$ tal que:
\begin{enumerate}
  \item $\|v\| \geq 0$, y $\|v\| = 0 \iff v = 0$.
  \item $\|\alpha v\| = |\alpha|\,\|v\|$ para todo escalar $\alpha$.
  \item Desigualdad triangular: $\|u+v\| \leq \|u\| + \|v\|$.
\end{enumerate}

\subsection{La familia de normas \texorpdfstring{$L^p$}{Lp} en \texorpdfstring{$\mathbb{R}^n$}{Rn}}

Para $p \geq 1$:
$$\|x\|_p = \left(\sum_{i=1}^n |x_i|^p\right)^{1/p}$$

\begin{itemize}
  \item $p=2$: norma \textbf{euclídea}.
  \item $p=1$: norma del \textbf{taxi} / \textbf{Manhattan}: $\|x\|_1 = \sum_i |x_i|$.
  \item $p=\infty$ (límite formal): norma \textbf{infinito} / \textbf{de
  Chebyshev}: $\|x\|_\infty = \max_i |x_i|$.
\end{itemize}

\begin{notebox}[Bolas unitarias]
En $\mathbb{R}^2$: $\|x\|_1=1$ es un rombo, $\|x\|_2=1$ es la circunferencia
usual, y a medida que $p$ crece las bolas se ``inflan'' hasta el cuadrado
$\|x\|_\infty=1$.
\end{notebox}

\section{Normas de matrices}

Las matrices cuadradas también forman un espacio vectorial, así que
admiten norma.

\subsection{Norma de Frobenius}

``Desatar'' la matriz en un vector de $n^2$ entradas y aplicar una norma
$L^p$ vectorial. Para $p=2$ se llama \textbf{norma de Frobenius}; para otros
$p$ no tiene uso conocido.

\begin{warnbox}[Limitación]
Para lo que busca esta unidad (relacionar residuo y error, condicionamiento),
las normas tipo Frobenius no son las útiles. El camino que importa es el de
las normas operador.
\end{warnbox}

\subsection{Normas operador (inducidas)}

Identificando $A$ con $x \mapsto Ax$, y dada una norma vectorial
$\|\cdot\|_B$, se define la \textbf{norma matricial inducida}:
$$\|A\|_M = \max_{x \neq 0} \frac{\|Ax\|_B}{\|x\|_B} = \max_{\|x\|_B \leq 1} \|Ax\|_B = \max_{\|x\|_B = 1} \|Ax\|_B$$

\begin{notebox}
No se parece a las normas $L^p$ ``ingenuas'': no combina entradas de $A$
directamente, sino que mide el máximo estiramiento que $A$ produce sobre un
vector unitario.
\end{notebox}

\section{La norma inducida por \texorpdfstring{$L^\infty$}{L-infinito}: máximo de las normas \texorpdfstring{$L^1$}{L1} de las filas}

\begin{keybox}[Proposición]
$$\|A\|_\infty = \max_{1\leq i\leq n} \sum_{j=1}^n |a_{ij}|$$
\end{keybox}

\subsection{Demostración}

Se prueban las dos desigualdades por separado.

\textbf{($\leq$)}: Sea $x$ con $\|x\|_\infty = 1$. Entonces
$$\|Ax\|_\infty = \max_i \left|\sum_{j=1}^n a_{ij}x_j\right| \leq \max_i \sum_{j=1}^n |a_{ij}||x_j| \leq \max_i \sum_{j=1}^n |a_{ij}|$$
usando la desigualdad triangular y $|x_j|\leq 1$. Tomando máximo sobre esos
$x$: $\|A\|_\infty \leq \max_i \sum_j |a_{ij}|$.

\textbf{($\geq$)}: Sea $i_0$ la fila que alcanza el máximo. Si $A=0$ no hay
nada que probar. Si $A\neq 0$, se construye
$$\hat{x}_j = \operatorname{sg}(a_{i_0 j}), \qquad j = 1,\dots,n$$
Como $\hat{x}_j \in \{-1,0,1\}$, $\|\hat{x}\|_\infty = 1$. Entonces:
$$\|A\hat{x}\|_\infty = \max_i \left|\sum_j a_{ij}\hat{x}_j\right| \geq \left|\sum_j a_{i_0 j}\hat{x}_j\right| = \sum_j |a_{i_0 j}| = \max_i \sum_j |a_{ij}|$$
donde se usó que $a_{i_0 j}\cdot\operatorname{sg}(a_{i_0 j}) = |a_{i_0 j}|$.
Como $\|A\|_\infty$ es el máximo sobre todos los $x$ unitarios, es
$\geq$ lo obtenido con $\hat{x}$.

\begin{notebox}[Técnica general]
Una desigualdad con máximo suele salir directo de la definición; la otra
exige exhibir un vector particular que alcanza la cota.
\end{notebox}

\section{Otras normas operador (enunciadas, no probadas en clase)}

\begin{itemize}
  \item \textbf{Inducida por $L^1$}: $\|A\|_1 = \max_j \sum_i |a_{ij}|$
  (máximo de las normas $L^1$ de las \textbf{columnas}) — corolario vía
  $\|A\|_\infty = \|A^T\|_1$. Demostración análoga, ejercicio.
  \item \textbf{Inducida por $L^2$}: $\|A\|_2$ es el \textbf{primer valor
  singular} de $A$. $A^TA$ es simétrica semidefinida positiva
  (diagonalizable, valores propios $\lambda_i \geq 0$); los valores
  singulares son $\sqrt{\lambda_i}$. Prueba en los apuntes, no desarrollada
  en clase (útil en compresión de datos e imágenes vía SVD, más adelante en
  el curso).
\end{itemize}

\section{Compatibilidad y submultiplicatividad}

Sea $\|\cdot\|_M$ una norma matricial inducida por $\|\cdot\|_B$.

\begin{keybox}[Compatibilidad]
Para toda $A$ y todo $x\in\mathbb{R}^n$:
$$\|Ax\|_B \leq \|A\|_M \|x\|_B$$
\end{keybox}

\textit{Demostración}: si $x=0$ es trivial. Si $x\neq 0$,
$\frac{\|Ax\|_B}{\|x\|_B} \leq \max_{y\neq 0} \frac{\|Ay\|_B}{\|y\|_B} = \|A\|_M$
por definición de máximo; multiplicando por $\|x\|_B$ se obtiene el
resultado.

\begin{keybox}[Submultiplicatividad]
Para toda $A,B$ matrices:
$$\|AB\|_M \leq \|A\|_M \|B\|_M$$
\end{keybox}

\textit{Demostración}: aplicando compatibilidad dos veces (viendo $Bx$ como
un vector al que se le aplica $A$):
$$\|AB\|_M = \max_{x\neq0} \frac{\|A(Bx)\|_B}{\|x\|_B} \leq \max_{x\neq0} \frac{\|A\|_M\|Bx\|_B}{\|x\|_B} \leq \max_{x\neq0} \frac{\|A\|_M\|B\|_M\|x\|_B}{\|x\|_B} = \|A\|_M\|B\|_M$$

\begin{notebox}[Corolario]
Para toda norma operador: $\|A^k\| \leq \|A\|^k$ (inducción, aplicando
submultiplicatividad con $B=A$ repetidamente).
\end{notebox}

\begin{warnbox}[Advertencia]
Estas dos propiedades valen para \textbf{cualquier norma operador}
(inducida por una norma vectorial) — no necesariamente para cualquier norma
matricial. No está establecido que la norma de Frobenius (§3.1) las cumpla.
\end{warnbox}

\section{Error vs. residuo: la motivación de fondo}

Al resolver $Ax=b$ con eliminación gaussiana con pivoteo parcial en una
computadora de precisión finita, se obtiene $\bar{x}$, distinta en general de
la solución exacta $x^\star$.

\begin{itemize}
  \item \textbf{Error}: $e = \bar{x} - x^\star$.
  \item \textbf{Residuo}: $r = A\bar{x} - b$.
\end{itemize}

\begin{warnbox}[El error no es computable]
Para calcularlo hace falta conocer $x^\star$, justo lo que no se tiene en un
problema real.
\end{warnbox}

\begin{keybox}[El residuo sí es computable]
Sólo requiere $A$, $b$ y la salida $\bar{x}$ del algoritmo.
\end{keybox}

\textbf{Relación}: como $b = Ax^\star$,
$$r = A\bar{x} - Ax^\star = A(\bar{x}-x^\star) = A\,e$$
$$\boxed{\text{el residuo es la matriz } A \text{ aplicada al error}}$$

Si $A$ es invertible, $r=0 \iff e=0$ (porque $e = A^{-1}r$): en aritmética
exacta, residuo nulo garantiza error nulo. En la práctica el residuo casi
nunca da exactamente cero.

\begin{notebox}[La pregunta que abre la clase siguiente]
Si el residuo tiene norma chica, ¿el error también es chico? \textbf{No hay
por qué} — no es obvio en absoluto. Es la pregunta que determina si un
método es ``bueno''. La Clase 8 (Número de Condición) da la respuesta usando
las herramientas de este capítulo.
\end{notebox}

\end{document}
